% Source: 26.singular value decomposition.v01, a two-page section in which
% Prof. Biran numbers nothing: the single theorem below carries no number in
% the notes. We deliberately print it as 26.a.1 anyway -- the chapter holds
% exactly one result, so numbering it shifts nothing, and it keeps the
% Singular Value Decomposition citable by \cref from later chapters.
\renewcommand{\thetheorem}{26.a.\arabic{theorem}}
\setcounter{theorem}{0}
\setcounter{chapter}{25}
% ---------------------------------------------------------------------------
% Provenance of the prose in this file.
%   % Extractor: <model>   the block below was transcribed from Prof. Biran's
%                          handwritten notes by that model
%   % Creator:   <model>   the block below was written by that model and has no
%                          counterpart in the notes
% A tag holds until the next tag.
% ---------------------------------------------------------------------------

% Extractor: Gemini 3.1 Pro
\chapter{Singular Value Decomposition}

% Thm. (26.a.1)
\begin{theorem}[Singular Value Decomposition]
\label{thm:singular_value_decomposition}
Let $K = \mathbb{R}$ or $\mathbb{C}$, and let $A \in \M_{m \times n}(K)$. Then there exist $P \in \Orth(m)$ (if $K = \mathbb{R}$) or $P \in \Unit(m)$ (if $K = \mathbb{C}$), and $Q \in \Orth(n)$ (if $K = \mathbb{R}$) or $Q \in \Unit(n)$ (if $K = \mathbb{C}$), and a \qt{diagonal} matrix
\[
D = \left(
\begin{array}{ccc}
\lambda_1 & & 0 \\
& \ddots & \\
0 & & \lambda_k \\
\hline
\multicolumn{3}{c}{0 \dots 0}
\end{array}
\right) \in \M_{m \times n}(\mathbb{R})
\]
with $\lambda_1, \dots, \lambda_k > 0$, where $k = \rank A$, such that $A = P D Q^{-1}$.
\end{theorem}

\begin{remark}
Recall that we previously defined the equivalence of matrices. Here, the factorization $A = P D Q^{-1}$ demonstrates that we have an orthogonal (or unitary) equivalence.
% Creator: Opus 5
The definition in question is \cref{def:equivalence_and_similarity} \textbf{(a)}: two matrices $A, B \in \M_{m \times n}(K)$ are equivalent when $B = PAQ$ for some $P \in \GL_m(K)$ and $Q \in \GL_n(K)$. The theorem sharpens this considerably, because the two invertible matrices may be taken to be orthogonal (resp. unitary) rather than merely invertible.
% Extractor: Gemini 3.1 Pro
\end{remark}

\begin{proof}
We will prove only a special case: we assume $\rank A = n$ (which is true if and only if $T_A$ is injective. Hence, we also assume $n \le m$). 

Define $B := A^* \cdot A \in \M_{n \times n}(K)$. We claim that the linear map $T_B \colon K^n \to K^n$ is injective.
Indeed, for all $0 \neq v \in K^n$, we have:
\begin{equation}
\label{eq:svd_pos_def}
\langle T_B v, v \rangle = \langle Bv, v \rangle = \langle A^* A v, v \rangle = \langle A v, A v \rangle > 0.
\end{equation}
This implies that $T_B v \neq 0$. Therefore, $\ker T_B = \{0\}$, hence $T_B$ is injective.

% Creator: Opus 5
Two steps in the chain \eqref{eq:svd_pos_def} deserve a word. The third equality is the defining property \eqref{eq:adjoint_property} of the adjoint, which is available for $T_A$ because the standard basis of $K^n$ is orthonormal, so that $(T_A)^* = T_{A^*}$ by \cref{prop:matrix_representation_adjoint}. The final inequality is strict precisely because of the standing assumption: $\rank A = n$ says that $T_A$ is injective by \cref{thm:rank_nullity}, so $v \neq 0$ forces $Av \neq 0$, and positive definiteness of the inner product then yields $\langle Av, Av \rangle > 0$.
% Extractor: Gemini 3.1 Pro

Note also that if $\eta \in K$ is an eigenvalue of $B$, then $\eta \in \mathbb{R}$ (even if $K = \mathbb{C}$) and $\eta > 0$.
Indeed, $B$ is self adjoint ($B^* = (A^* A)^* = A^* (A^*)^* = A^* A = B$), hence $\eta \in \mathbb{R}$. (This last step is \cref{lem:self_adj_properties} \textbf{(a)}, applied to the self-adjoint endomorphism $T_B$ of $K^n$). To see that $\eta > 0$, let $v$ be an eigenvector for $\eta$, so that $v \neq 0$ and $Bv = \eta v$. Substituting this into \eqref{eq:svd_pos_def} gives
\[
0 < \langle Bv, v \rangle = \langle \eta v, v \rangle = \eta \cdot \langle v, v \rangle,
\]
where the last step uses linearity in the first variable together with $\eta \in \mathbb{R}$. Since $v \neq 0$, positivity of the inner product gives $\langle v, v \rangle > 0$, and therefore, $\eta > 0$.

By the Spectral Theorems for matrices (\cref{thm:spectral_thm_matrices}), $B$ is orthogonally diagonalizable, hence there exists $Q \in \Orth(n)$ (or $\Unit(n)$) such that
\[
Q^{-1} A^* A Q = \begin{pmatrix} \eta_1 & & 0 \\ & \ddots & \\ 0 & & \eta_n \end{pmatrix},
\]
where $\eta_j \in \mathbb{R}$, and $\eta_j > 0$ for all $j$.
% Creator: Opus 5
(For $K = \mathbb{R}$ this is part \textbf{(b)} of that theorem, because $B = A\transp A$ is then a real symmetric matrix; for $K = \mathbb{C}$ it is part \textbf{(a)}, because a self-adjoint matrix is in particular normal, $B B^* = B B = B^* B$).
% Extractor: Gemini 3.1 Pro

Let $\lambda_j := +\sqrt{\eta_j} > 0$ for $j = 1, \dots, n$.
Define
\[
C := A Q \cdot \begin{pmatrix} \lambda_1^{-1} & & 0 \\ & \ddots & \\ 0 & & \lambda_n^{-1} \end{pmatrix} \in \M_{m \times n}(K).
\]
We have:
\begin{align*}
C^* C &= \begin{pmatrix} \lambda_1^{-1} & & 0 \\ & \ddots & \\ 0 & & \lambda_n^{-1} \end{pmatrix} \underbrace{Q^{-1} A^* A Q}_{=\begin{pmatrix} \eta_1 & & 0 \\ & \ddots & \\ 0 & & \eta_n \end{pmatrix}} \begin{pmatrix} \lambda_1^{-1} & & 0 \\ & \ddots & \\ 0 & & \lambda_n^{-1} \end{pmatrix} \\
&= I_n.
\end{align*}

It follows that the columns of $C$ form an orthonormal system of vectors in $K^m$. To see this explicitly, let $c_1, \dots, c_n$ be the columns of $C$. The identity $C^* C = I_n$ means:
\[
C^* C = 
\begin{pmatrix}
- & \overline{c_1}\transp & - \\
& \vdots & \\
- & \overline{c_n}\transp & -
\end{pmatrix}
\begin{pmatrix}
| & & | \\
c_1 & \dots & c_n \\
| & & |
\end{pmatrix}
= I_n.
\]
% Creator: Opus 5
Reading off the entries of this product, the identity $C^* C = I_n$ says precisely that $\langle c_j, c_i \rangle = \overline{c_i}\transp c_j = \delta_{ij}$ for all $1 \leq i, j \leq n$, which is the assertion that $c_1, \dots, c_n$ is an orthonormal system.
% Extractor: Gemini 3.1 Pro
Extend $c_1, \dots, c_n$ to an orthonormal basis $c_1, \dots, c_n, f_{n+1}, \dots, f_m$ of $K^m$ and define:
\[
P := \begin{pmatrix} | & & | & | & & | \\ c_1 & \dots & c_n & f_{n+1} & \dots & f_m \\ | & & | & | & & | \end{pmatrix} \in \M_{m \times m}(K).
\]
Clearly, $P \in \Orth(m)$ (or $\Unit(m)$).
% Creator: Opus 5
(The extension is possible by \cref{exc:extend_orthonormal_system}, and $P$ is orthogonal, resp. unitary, directly by \cref{def:orthogonal_unitary_matrices}, since its columns are by construction an orthonormal basis of $K^m$).
% Extractor: Gemini 3.1 Pro

We have:
\[
A Q \begin{pmatrix} \lambda_1^{-1} & & 0 \\ & \ddots & \\ 0 & & \lambda_n^{-1} \end{pmatrix} = C = P \cdot 
\left(
\begin{array}{ccc}
1 & & 0 \\
& \ddots & \\
0 & & 1 \\
\hline
\multicolumn{3}{c}{0}
\end{array}
\right)
\begin{array}{l}
\left. \vphantom{\begin{matrix} 1 \\ \ddots \\ 1 \end{matrix}} \right\} n \\
\left. \vphantom{\begin{matrix} 0 \end{matrix}} \right\} m-n
\end{array}
\]
Multiplying from the right by the diagonal matrix of $\lambda_j$'s and then by $Q^{-1}$, we obtain:
\[
A = P \cdot \left(
\begin{array}{ccc}
\lambda_1 & & 0 \\
& \ddots & \\
0 & & \lambda_n \\
\hline
\multicolumn{3}{c}{0}
\end{array}
\right) \cdot Q^{-1}.
\]
This proves the theorem (under the assumption that $\rank A = n$).

For the general case, see the script, chapter 7.5.
\end{proof}

